\chapter{Control-Flow Graphs}
\label{ch:control-flow-graphs}

A flat instruction list shows every operation but hides where execution can go next.
Most instructions fall through to their successor, yet branches and jumps redirect control to distant labels.
This chapter groups instructions into \term{basic blocks} and connects the blocks with edges into a \term{control-flow graph}.
Almost every analysis in the rest of the book runs on this graph rather than on raw lists.
Getting its construction right pays off in every chapter that follows.

\section{Straight runs and jumps}
\label{sec:ch07-straight-runs}

Consider the lowered fragment below, where \code{t1} holds a test value.
\begin{irlisting}
label L1;
t1 = x > 0;
br t1, L2, L3;
label L2;
y = 1;
jmp L4;
label L3;
y = 2;
label L4;
print y
\end{irlisting}
Execution cannot wander arbitrarily through these $10$ instructions.
It enters at \code{L1}, evaluates the test, and jumps to exactly one of \code{L2} or \code{L3}.
Each arm assigns \code{y} once and rejoins at \code{L4}, which prints.
The straight runs are the maximal instruction sequences with a single entry and a single exit.
Here they are the run from \code{L1} through the branch, the run \code{y = 1} plus its jump, the run \code{y = 2}, and the run \code{print y}.
Each run becomes one basic block: a sequence that always executes from its first instruction to its last.
A block ends with a branch, a jump, or the instruction before a label, and a block begins after a branch, after a jump, or at a label.
Table~\ref{tab:ch07-blocks} lists the four blocks of the fragment with their contents and successors.
\begin{table}
\caption{Basic blocks of the branching fragment with their instructions and successor blocks.}
\label{tab:ch07-blocks}
\centering
\begin{tabular}{lll}
\toprule
Block & Instructions & Successors \\
\midrule
B1 & test and conditional branch & B2, B3 \\
B2 & assign \code{y = 1} and jump & B4 \\
B3 & assign \code{y = 2}, falls through & B4 \\
B4 & print \code{y} & none, exit \\
\bottomrule
\end{tabular}
\end{table}
The table makes the join explicit: both arms list \code{B4} as a successor.
Falling through counts as an edge too, so \code{B3} reaches \code{B4} without any jump instruction.
The graph therefore has $4$ nodes and $4$ edges for $10$ instructions.
Every execution corresponds to one path from the entry block to the exit, and every graph path that respects branch conditions corresponds to an execution.
\section{Predecessors and successors}
\label{sec:ch07-pred-succ}

Each block records two sets: its \term{successors}, where control may go next, and its \term{predecessors}, where control may come from.
We write them with the shared notation $\Succ(B)$ and $\Pred(B)$.
For the fragment above, $\Succ(B1) = \{B2, B3\}$ and $\Pred(B4) = \{B2, B3\}$.
Successors come directly from terminators: a conditional branch contributes two, a jump contributes one, and a fall-through contributes the next block.
Predecessors are the inverse relation, computed by walking every successor edge once and recording it backward.
Both directions are needed because analyses propagate facts along opposite directions.
A forward analysis follows successors from the entry, while a backward analysis follows predecessors from the exit.
Every graph in this book has a single entry block with no predecessors and a single exit block with no successors.
Lowering guarantees the single entry by starting each function body in its own fresh block.
It guarantees the single exit by routing every return and fall-off-the-end path through one final block.
These conventions simplify every algorithm that follows: iteration can start from a known point and stop at a known point.
A program with an infinite loop still has an exit block in the graph even when no execution ever reaches it.
The exit block then has predecessors from the loop but contributes no dynamic executions, which is exactly the shape a liveness analysis expects.
Keeping unreachable exits in the graph also preserves positions for warnings about missing returns.
An analysis that needs reachable-only views computes reachability itself rather than asking the builder to delete blocks.
Figure~\ref{fig:ch07-diamond-loop} shows the diamond of the fragment extended with a loop back to the test.
\begin{figure}
\centering
\begin{tikzpicture}[node distance=1.7cm]
\node[basic block] (b1) {B1: test};
\node[basic block, below left of=b1, xshift=-1.2cm] (b2) {B2: y = 1};
\node[basic block, below right of=b1, xshift=1.2cm] (b3) {B3: y = 2};
\node[basic block, below of=b2, xshift=1.2cm] (b4) {B4: print};
\draw[cfg edge] (b1) -- (b2);
\draw[cfg edge] (b1) -- (b3);
\draw[cfg edge] (b2) -- (b4);
\draw[cfg edge] (b3) -- (b4);
\draw[back edge] (b4) -- (b1);
\end{tikzpicture}
\caption{The diamond test with both arms joining at the print block and a dashed edge looping back to the test.}
\label{fig:ch07-diamond-loop}
\end{figure}
The four solid edges are the diamond from Table~\ref{tab:ch07-blocks}.
The dashed edge returns from the print block to the test, turning the fragment into a loop that retests after each print.
Edges that return to an earlier block are called \term{back edges} and are drawn dashed throughout this book.
The next chapter uses them to identify loops precisely.

\section{Building the graph}
\label{sec:ch07-building-graph}

Algorithm~\ref{alg:ch07-cfg} constructs the graph from a flat instruction list using the leader method.
\begin{algorithm}
\caption{Leader-based construction of the control-flow graph.}
\label{alg:ch07-cfg}
\begin{pseudocode}
\Require flat instruction list $C$ with labels, branches, and jumps
\Ensure block list with $\Pred$ and $\Succ$ sets for every block
\Function{BuildCFG}{$C$}
\State mark the first instruction and every labeled instruction as a leader
\State mark every instruction after a terminator, including a return, as a leader
\State cut the list into blocks starting at each leader
\State link each block to its successors from its terminator or fall-through
\State invert the successor links to obtain all predecessor sets
\EndFunction
\end{pseudocode}
\end{algorithm}
The algorithm assumes every branch and jump targets a label that exists and that the list is nonempty.
It guarantees one block per maximal straight run and exactly one edge per possible transfer of control.
The leader marks are the whole insight: the first instruction starts a block, every label starts a block, and every instruction after a transfer starts a block.
Cutting at leaders therefore produces the longest runs with single entries and single exits.
Linking reads only block terminators, so each block is processed once and each edge is created once.
Inversion then visits every edge once more to fill in predecessors.
Handle an empty program separately by creating its designated empty entry/exit representation.
Unreachable blocks can have predecessors, including an entire unreachable cycle; reachability requires traversal from entry, not an in-degree test.
Unreachable blocks are kept rather than deleted at this stage, because a later pass may warn about them with exact positions.

\section{Worked example}
\label{sec:ch07-worked-example}

We now build the graph for a small loop program by hand.
\begin{irlisting}
i = 1;
label head;
t1 = i < 5;
br t1, body, done;
label body;
i = i + 1;
jmp head;
label done;
print i
\end{irlisting}
The list has $9$ instructions.
Leaders are instruction $1$ (the first), instruction $2$ (a label target), instruction $5$ (a label target), instruction $7$ (the instruction after the jump on line $6$ counts, but line $7$ is also a label target), and instruction $8$ is part of the same block as line $7$.
Cutting gives $4$ blocks: $B1$ holds line $1$, $B2$ holds lines $2$ through $4$, $B3$ holds lines $5$ and $6$, and $B4$ holds lines $7$ through $9$ after splitting at the label, which is $2$ lines.
More carefully, $B4$ holds \code{label done} and \code{print i}, which is $2$ instructions.
Successors are $\Succ(B1) = \{B2\}$ by fall-through, $\Succ(B2) = \{B3, B4\}$ by branch, $\Succ(B3) = \{B2\}$ by jump, and $\Succ(B4) = \emptyset$ at exit.
Inverting gives $\Pred(B1) = \emptyset$, $\Pred(B2) = \{B1, B3\}$, $\Pred(B3) = \{B2\}$, and $\Pred(B4) = \{B2\}$.
The edge count is $4$: one fall-through, two branch arms, and one jump.
Running the program confirms the shape: \code{i} takes values $1$ through $5$, the body runs $4$ times, and the back edge $B3$ to $B2$ is taken $4$ times while the exit edge $B2$ to $B4$ is taken once.
The hand computation matches Figure~\ref{fig:ch07-diamond-loop} in structure: a test block with two successors, a body block with a back edge, and an exit block.

\section{Implementation notes}
\label{sec:ch07-implementation}

Blocks are usually objects holding an instruction list plus adjacency sets for successors and predecessors.
Instructions stay in arrays inside blocks for compact sweeping, while edges are stored as pointer or index pairs between blocks.
Keeping both $\Pred$ and $\Succ$ doubles the edge memory but makes forward and backward passes equally cheap.
The running time is $O(N + E)$, where $N$ is the number of instructions and $E$ is the number of control edges.
Leader marking scans all $N$ instructions once, cutting walks the marks once, and linking plus inversion each visit every edge once.
Memory use is $O(N + E)$ for the blocks, instructions, and both adjacency relations.
A subtle case is the \term{critical edge}: an edge from a block with several successors to a block with several predecessors.
Splitting such edges by inserting an empty block simplifies many later transformations, and production builders offer it as an option.
Another subtlety is abnormal control such as errors and early exits, which Mini omits.
Real languages add edges for exceptions, short-circuit operators, and indirect jumps, and each needs a documented policy in the builder.
Unreachable-block elimination and jump threading are tempting at construction time but belong in later passes.
The builder should preserve the program shape faithfully so that warnings and listings map back to source positions.
\begin{ccnote}[Fall-through is an edge]
A block that ends without a jump still transfers control to the next block.
Forgetting that edge silently drops half of every diamond and makes the join block look unreachable.
\end{ccnote}

\section{Exercises}
\label{sec:ch07-exercises}

\begin{enumerate}
\item \label{ex:ch07-build-diamond} Build the graph for the fragment of Section~\ref{sec:ch07-straight-runs} using Algorithm~\ref{alg:ch07-cfg}. List the leaders, the blocks, and the $\Pred$ and $\Succ$ sets. Confirm $4$ blocks and $4$ edges and check your sets against Table~\ref{tab:ch07-blocks}.
\item \label{ex:ch07-count-paths} Add the back edge of Figure~\ref{fig:ch07-diamond-loop} to the diamond and count the distinct paths from entry to the print block that use the back edge $0$, $1$, and $2$ times. State the general formula for $K$ uses and explain why the count grows.
\item \label{ex:ch07-critical} Give a three-block Mini program whose graph contains a critical edge. Identify the block with several successors and the block with several predecessors. Show where the inserted empty block goes and how the $\Pred$ and $\Succ$ sets change.
\end{enumerate}

\section{Further reading}
\label{sec:ch07-further-reading}

Appel constructs control-flow graphs as the foundation for dataflow analysis and register allocation, which is the same role they play in the coming chapters, see~\cite{appel-1998-modern}.
Morgan treats graph construction as production infrastructure that must handle every control form the language allows, see~\cite{morgan-1998-optimizing}.
Safonov shows how a trustworthy compiler validates its own intermediate structures rather than assuming each pass built them correctly, see~\cite{safonov-2010-trustworthy}.
The C language definition specifies the control constructs whose lowering produces the branches, jumps, and labels of this chapter, see~\cite{wg14-2011-c11}.
